\section{R tmle: preserve the target and audit uncertainty before reuse}
\label{sec:exec24-tmle-reuse}

A banking team considering a single repayment offer may need the change in
repayment probability, the ratio of repayment risks, or the effect among
customers who actually received the offer. These targets can lead to different
allocation decisions. The selected R \texttt{tmle} 2.1.1 distribution supplies
binary point-treatment machinery, outcome-missingness adjustment, ATT and ATC
paths, and a binary intermediate-variable route. It is a useful reference for
an eventual native R worker. It does not, merely by accepting repeated customer
identifiers, establish identification for a sequence of adaptive interventions.
Our earlier inspection of missing treatment and baseline covariates remains
applicable: these are rejected, and the no-variation remapping is a mean-outcome
route rather than reconstruction of unknown assignments.
\citep{exec20_rfx_tmle}

The selected CRAN mirror revision is dated 30 August 2025 and precedes the
literature cutoff. We inspect its actual treatment function, help text,
dependency declaration, rights file and citation metadata. The latter identifies
a 2012 package article, but is not a substitute for reading its original.
The comparison here rests on the pinned technical implementation and authored
arithmetic; it does not borrow uninspected results from the package bibliography.
\citep{exec24_tmle_sources}

\subsection{Point-treatment meaning, weighting and the population clock}
Let $W$ contain information available before the offer, $A\in\{0,1\}$ its
assignment, $Y$ the later repayment indicator, and $\Delta$ its observation
indicator. Under consistency, conditional exchangeability, treatment positivity
and an appropriate outcome-observation assumption, an identified point-treatment
mean may be represented by an average of $Q_a(W)$, where $Q_a$ is the conditional
outcome regression at treatment $a$. These are substantive premises for the
proposal, not conclusions of a successful regression. Sampling weights $w_i$
produce the empirical plug-in targets
\[
 \mu_a=\frac{\sum_i w_i Q_a(W_i)}{\sum_i w_i},\qquad
 \mathrm{RD}=\mu_1-\mu_0,\quad
 \mathrm{RR}=\frac{\mu_1}{\mu_0},\quad
 \mathrm{OR}=\frac{\mu_1/(1-\mu_1)}{\mu_0/(1-\mu_0)}.
\]
The source normalizes observation weights to sum to the number of rows and
uses corresponding weighted means. Its parameter helper calculates these
contrasts and distinct influence-function expressions. The adapter must report
the actual target, denominator, observation rule and weight interpretation.
\citep{exec24_tmle_sources}

In an authored two-profile example, weights $(1,3)$ and regression predictions
$Q_0=(0.1,0.4)$, $Q_1=(0.2,0.5)$ give $\mu_0=13/40$, $\mu_1=17/40$,
$\mathrm{RD}=1/10$, $\mathrm{RR}=17/13$ and $\mathrm{OR}=459/299$.
At a stipulated 200 USD per incremental repayment, the difference gives
20 USD per offer before costs; the ratio itself is not a dollar benefit.
Neither a fitted conditional mean nor an externally supplied prediction matrix
establishes the causal assumptions. A later prototype must preserve training,
fold, prediction and customer identities and qualify the actual nuisance
estimation and uncertainty route.
\citep{exec24_tmle_contract}

Observation and customer populations also differ. One customer with one row
and effect zero, and another with three rows each having effect ten, give a
row-average effect of $7.5$ and an equal-customer average of $5$. The inspected
helper averages influence-function contributions within identifiers and divides
the variance of those means by the number of identifiers. Whether that aligns
with a weighted row target requires its sampling law and cluster-size argument;
an identifier field alone is insufficient.
For row scores $(-3,1,1,1)$ in clusters of sizes one and three, the cluster means
are $(-3,1)$ and their sample variance divided by two is $4$. For comparison,
the cluster-sum contributions scaled by the observed cluster-to-row ratio are
$(-1.5,1.5)$, whose sample variance divided by two is $9/4$. This is exact
arithmetic illustrating different contributions, not a proof of native
miscalibration or a universal prescription for dependent observations.
\citep{exec24_tmle_sources,exec24_tmle_contract}

\subsection{Intermediate-variable targets and a specific interval discrepancy}
When the supplied intermediate variable has more than one observed value, the
source constructs intervention predictions at its two binary levels and uses
joint treatment, intermediate-variable and outcome-observation probabilities.
A controlled intervention on a mediator needs a defensible action and the
corresponding sequential exchangeability and support assumptions. Conditioning
on an observed mediator value is not automatically the same intervention.
Truncating a probability of $0.01$ to $0.1$ changes an inverse contribution from
100 to 10; it does not create unobserved interventions or prove identification.
The source's explicit bounds and retained rows belong in the result contract.
\citep{exec24_tmle_sources,exec24_tmle_contract}

A concrete static discrepancy matters for this route. Its bootstrap matrix is
named, in order, \texttt{EY0}, \texttt{EY1}, \texttt{ATE}, \texttt{logRR} and
\texttt{logOR}. The source fills the fourth column with log risk ratios and
uses column four for the risk-ratio variance, but constructs its risk-ratio
bootstrap interval by exponentiating quantiles of column three. The ordinary
point-treatment branch uses column four for this interval. We verified the
column definitions, assignments, branch condition and output statements; this
is a mismatch in the selected source, not an observed package execution.
\citep{exec24_tmle_sources}

The difference is visible without running that implementation. If every row of
an authored bootstrap fixture has risks $0.1$ and $0.2$, its risk difference is
$0.1$ and its risk ratio is $2$. Exponentiating the constant difference gives
$\exp(0.1)\simeq1.10517$ instead of $2$. Thus column substitution changes the
scientific quantity even when numerical routines return finite values. The
proposed integration refuses the affected interval pending a separately
reviewed correction or qualified successor. A repair must use named quantities,
check both intermediate-variable levels and all interval fields, and preserve
the unmodified original with the evidence. We do not patch upstream code or
silently present our arithmetic as its native output.
\citep{exec24_tmle_contract}

\subsection{Bootstrap counts, missing outputs and uncertainty admission}
Both inspected branches iterate from one through $B$. Iterations below $B$
resample identifiers, while the last uses the original sample so returned fits
have the original identities. Quantiles are calculated over the full matrix.
Consequently the inspected loop contains $B-1$ resampled rows and one original
row, which must be distinguished from a request for $B$ independent resamples.
For an authored vector $(0,1,100)$, where the last value represents the original
sample, linear-interpolated 2.5 and 97.5 percent quantiles are $0.05$ and $95.05$.
Using only the two stipulated resamples gives $0.025$ and $0.975$. These numbers
isolate a counting convention; they do not measure empirical coverage.
\citep{exec24_tmle_sources,exec24_tmle_contract}

The inspected variance call does not remove missing values, while its quantile
calls explicitly do. With the illustrative vector $(1,\mathrm{NA},3)$ the
variance is missing but the finite-only quantile interval is $[1.05,2.95]$.
A successful interval field cannot hide its incomplete denominator. The future
worker must return requested, attempted, original-sample, finite and failed
counts, and identify the actual nuisance-fitting and targeting steps repeated.
The intermediate-variable bootstrap samples identifiers separately from the two
observed mediator groups and then includes all rows for each sampled identifier.
Identifiers spanning both groups require explicit scrutiny; disjoint resampling
populations cannot be presumed from a row-level mediator label.
\citep{exec24_tmle_sources,exec24_tmle_contract}

The proposal therefore separates influence-function intervals, targeted
bootstrap intervals, sensitivity scenarios and identification failures. Boundary
risks, zero ratio denominators, unsupported missingness, nuisance failures,
truncation, nonconvergence and incomplete replicates need typed outcomes.
Neither an interval label nor the documentation's suggested high iteration
count proves coverage for arbitrary bank panels or adaptive policies.
\citep{exec24_tmle_contract}

\subsection{Rights, alternatives and a bounded adoption decision}
The DESCRIPTION declares \texttt{BSD\_3\_clause + file LICENSE | GPL-2}; the
selected LICENSE file contains year and copyright-holder fields. The treatment
source footer separately describes educational, research and not-for-profit
permission and directs commercial inquiries elsewhere. We preserve these
texts as a rights question for the selected distribution, without deciding their
legal precedence. No commercial permission is inferred from a package index.
Declared dependencies include \texttt{glmnet} and \texttt{SuperLearner}, with
optional learners and diagnostics; an installed dependency closure, its rights,
platform behavior and maintenance commitments remain to be qualified.
\citep{exec24_tmle_sources}

Our current disposition is reference. Conditional integration through the
original R worker is preferable to assuming that a whole-library port preserves
fitted objects, weighting and targeting semantics. A narrow independent
finite-population or analytical oracle can establish target and refusal
expectations. A separately qualified longitudinal estimator answers a different
request; it is not an alias for this point-treatment interface. Whole-library
porting is deferred until demonstrated demand and costs justify it. The
scientific owner must approve target, sampling and missingness assumptions; the
adapter owner must preserve identities and typed failure; the rights/platform
owner must resolve distribution and environment; and the evaluation owner must
preregister public positive, adverse, boundary and fault fixtures.
\citep{exec24_tmle_contract}

Illustrative total cost of 2,400 USD over 120 requests, 100 conformant results
and 80 useful results is respectively 20, 24 and 30 USD per denominator. These
are authored inputs, not measured prices or a business case. Rights review,
engineering, dependencies, compute including failures, storage, recovery and
maintenance belong in the same cost record. Public or generated experiments
can justify the next investment before expensive bank data work. They cannot
supply unknown assignment truth, universal hidden-confounding identification
or arbitrary dependent-data coverage. The interval discrepancy has a concrete
repair and validation route; scientific identification and commercial rights
remain explicit conditions, not promises of automatic closure.
\citep{exec24_tmle_contract}
